\section{Discussion}\label{discussion}

\subsection{6.1 Interpretation}\label{interpretation}

The central empirical fact is a panel-wide tail shift. A static model
can preserve reasonable rankings across items and locations while its
absolute q90 level becomes severely understated. This explains why a
simple common-shift adjustment can outperform a more flexible tree model
after a regime change. Tree ensembles interpolate historical outcomes
effectively but do not automatically extrapolate a new common tail level
beyond the training regime.

Different adaptations solve different problems. The fixed gate reacts
rapidly to a common shock; global conformal calibration targets overall
residual quantiles; hierarchy addresses cross-sectional heterogeneity;
and the static expert protects against needless adaptation in stable
periods. The controller reduces regret by diversifying these failure
modes. Its limited gain over the best expert reflects the small number
of monthly observations and the difficulty of identifying the current
regime before delayed losses mature.

For operational use, the model should be treated as a risk-monitoring
device. An elevated q90 forecast signals that a large three-month price
increase is plausible relative to the historical distribution. It does
not identify the causal source of the shock, and it should not replace
item-specific market intelligence. Forecasts should be accompanied by
coverage diagnostics, model weights, and an explicit recovery flag.

\subsection{6.2 Policy and risk-management
relevance}\label{policy-and-risk-management-relevance}

The framework is relevant to food-security monitoring because it
highlights the upper tail rather than only the mean. A central agency
could use it to identify months in which procurement budgets,
cash-transfer adequacy, or household food-access risks are unusually
exposed. Insurers, banks, and firms could use the same logic for stress
testing inventories, working capital, and commodity-linked liabilities.

The delayed-feedback design is also broadly relevant. Many risk systems
observe outcomes after a reporting lag, settlement period, or forecast
horizon. A method that quietly uses not-yet-matured errors will appear
more responsive than a deployable model. Enforcing maturation dates is
therefore a substantive methodological requirement, not merely a coding
detail.

\subsection{6.3 Limitations}\label{limitations}

\begin{itemize}
\item
  The Phase 10 post-2022 evaluation is exploratory because the
  architecture was motivated after earlier models had already been
  examined on the same period.
\item
  The NBS continuation is a same-source national-item validation rather
  than a replication of the original state-commodity panel.
\item
  The WFP external sample is concentrated in Borno and Yobe and is not
  nationally representative.
\item
  Exact publication and release lags are not reconstructed; the design
  assumes month-t predictors are available at issuance.
\item
  Several NBS item definitions change in later publications, and only
  sufficiently comparable series are retained.
\item
  The stress detector is standardised separately within each source and
  does not transport as a source-invariant economic index.
\item
  The convex combination of q90 forecasts is an operational tail
  forecast, not a conformal set with distribution-free conditional
  coverage.
\item
  Month-block bootstrap intervals preserve common within-month shocks
  but do not eliminate all serial dependence.
\item
  The monthly oracle is infeasible and is used only as a regret
  benchmark.
\item
  Simulation conclusions depend on the specified heavy-tailed panel
  processes and do not constitute a universal regret theorem.
\end{itemize}

\subsection{6.4 What would constitute confirmatory
evidence?}\label{what-would-constitute-confirmatory-evidence}

The controller should now be frozen. Confirmatory evidence requires an
untouched country or a future Nigerian evaluation period that was not
used to motivate, select, or refine any component. The primary outcome,
exclusions, mapping rules, target horizon, feedback delay, expert
definitions, controller parameters, bootstrap procedure, and superiority
criterion must be specified before outcomes are inspected. A separate
frozen protocol accompanies this manuscript.
